\rtmtight
\begin{tblr}{width=\textwidth,colspec={|Q[c,m,2.1cm]|X[l,m]|},hlines,vlines,hspan=minimal,rowsep=1.5pt,colsep=3pt,row{1}={bg=headgray,font=\bfseries}}
\SetCell{c,m}\hfil\logo\hfil & \textbf{POLITEKNIK NEGERI MALANG}\par\textbf{JURUSAN TEKNOLOGI INFORMASI}\par\textbf{PROGRAM STUDI: D4 TEKNIK INFORMATIKA} \\
\SetCell[c=2]{l} \textbf{RENCANA TUGAS MAHASISWA (RTM) DAN RUBRIK PENILAIAN} & \\
Nama Mata Kuliah & Keselamatan dan Kesehatan Kerja \\
Kode Mata Kuliah & RTI251008 \quad \textbf{sks / Jam perminggu:} 2 SKS/4 jam \quad \textbf{Semester:} 1 \\
Dosen Pengampu & \begin{enumerate}\item Anugrah Nur Rahmanto, S.Sn., M.Ds.\item Budi Harijanto, S.T., M.MKom.\item Sofyan Noor Arief, S.ST., M.Kom.\item Ariadi Retno Ririd, S.Kom., M.Kom.\end{enumerate} \\
\end{tblr}
\assessmentblock{Tugas Laporan dan Artefak K3}{Tugas terstruktur (laporan, makalah, slide animasi, role-play)}{SCPMK0101-00801, SCPMK0101-00802, SCPMK0101-00803, SCPMK0101-00804}{Mahasiswa menyelesaikan rangkaian tugas mingguan (Tugas-1 s.d. Tugas-11): laporan sejarah dan landasan hukum K3, slide animasi lingkungan kerja dan APD, demonstrasi P3K, makalah kesehatan masyarakat, dan struktur organisasi K3.}{Mengkaji materi dan peraturan K3, menyusun laporan/makalah/slide animasi, mendemonstrasikan atau memerankan skenario K3, dan mengumpulkan evidence melalui LMS.}{Laporan tugas, makalah, slide animasi, dan dokumentasi demonstrasi/role-play.}{Ketepatan konsep, sejarah, dan landasan hukum K3 & 6 & Konsep, sejarah, tujuan, lambang, serta undang-undang dan peraturan K3 diuraikan tepat dengan rujukan peraturan yang benar. & Konsep dan landasan hukum umumnya benar, tetapi sebagian rujukan peraturan kurang tepat atau tidak lengkap. & Konsep atau landasan hukum banyak yang keliru dan tidak merujuk peraturan yang berlaku. \\ \\
Ketepatan analisis bahaya, kecelakaan, dan SMK3 & 6 & Faktor penyebab kecelakaan, sumber bahaya, dan elemen SMK3 diidentifikasi lengkap dan dikaitkan dengan kasus industri. & Identifikasi bahaya dan elemen SMK3 cukup benar, tetapi kaitan dengan kasus industri masih terbatas. & Identifikasi bahaya atau elemen SMK3 keliru dan tidak terkait aktivitas industri. \\ \\
Kualitas artefak tugas (laporan, makalah, slide, role-play) & 6 & Artefak lengkap, sistematis, sesuai materi minggu berjalan, dan demonstrasi/role-play sesuai prosedur K3. & Artefak tersedia dan cukup sesuai materi, tetapi sistematika atau kesesuaian prosedur masih kurang. & Artefak tidak lengkap, tidak sesuai materi, atau prosedur K3 tidak diikuti. \\ \\
Ketepatan rancangan program dan dokumentasi K3 & 6 & Rancangan program kerja, struktur organisasi, tupoksi, dan dokumentasi K3 disusun runtut dan dapat diterapkan. & Rancangan program dan dokumentasi tersedia, tetapi tupoksi atau kelengkapan dokumen masih kurang. & Rancangan program tidak operasional atau dokumentasi K3 tidak tersedia. \\}

\assessmentblock{Kuis 1 Konsep dan Regulasi K3}{Kuis}{SCPMK0101-00801}{Mahasiswa mengerjakan kuis tertulis daring tentang konsep, sejarah, tujuan, lambang K3, serta undang-undang dan landasan hukum K3.}{Mengerjakan tes tulis daring secara mandiri dalam waktu yang ditentukan.}{Jawaban kuis daring.}{Penguasaan konsep, sejarah, dan tujuan K3 & 4 & Jawaban menjelaskan konsep, sejarah, tujuan, dan lambang K3 secara tepat dan lengkap. & Sebagian besar jawaban benar, tetapi penjelasan konsep atau sejarah belum lengkap. & Jawaban banyak yang keliru atau hanya menghafal istilah tanpa pemahaman. \\ \\
Ketepatan regulasi dan landasan hukum K3 & 4 & Undang-undang dan peraturan K3 disebutkan tepat dan dikaitkan dengan kewajiban di tempat kerja. & Regulasi utama disebutkan benar, tetapi kaitan dengan penerapannya belum jelas. & Regulasi yang disebut keliru atau tidak relevan dengan K3. \\}

\assessmentblock{UTS Konsep, Regulasi, dan Sistem Manajemen K3}{UTS (ujian tulis daring)}{SCPMK0101-00801, SCPMK0101-00802}{Mahasiswa mengerjakan ujian tengah semester yang mencakup konsep dan regulasi K3, analisis faktor penyebab kecelakaan, serta sistem manajemen K3 (SDM, operasi, dokumentasi).}{Mengerjakan ujian tulis daring secara mandiri, menjawab soal konsep dan soal analisis kasus.}{Lembar jawaban ujian daring.}{Penguasaan konsep dan regulasi K3 & 8 & Konsep, tujuan, dan landasan hukum K3 dijawab tepat serta diterapkan pada situasi kerja yang diberikan. & Konsep dan regulasi umumnya benar, tetapi penerapan pada situasi kerja belum tepat. & Konsep atau regulasi banyak yang keliru. \\ \\
Ketepatan analisis kecelakaan dan sistem manajemen K3 & 7 & Hubungan faktor penyebab kecelakaan dan keadaan tidak selamat dianalisis tepat serta elemen SMK3 diuraikan lengkap. & Analisis kecelakaan cukup tepat, tetapi uraian elemen SMK3 belum lengkap. & Analisis kecelakaan keliru atau elemen SMK3 tidak diuraikan. \\ \\
Kualitas jawaban studi kasus K3 & 5 & Solusi kasus K3 runtut, sesuai prosedur, dan mempertimbangkan pengelolaan SDM dan operasi K3. & Solusi kasus cukup relevan, tetapi langkah penyelesaian belum runtut. & Solusi kasus tidak relevan atau tidak menjawab permasalahan K3. \\}

\assessmentblock{Kuis 2 Pengendalian Risiko dan Organisasi K3}{Kuis}{SCPMK0101-00803, SCPMK0101-00804}{Mahasiswa mengerjakan kuis tertulis daring tentang manajemen risiko K3, alat pelindung diri, serta pengelolaan komunikasi, personel, dan organisasi K3 menuju zero accident.}{Mengerjakan tes tulis daring secara mandiri dalam waktu yang ditentukan.}{Jawaban kuis daring.}{Ketepatan pengendalian risiko dan APD & 4 & Langkah manajemen risiko dan pemilihan APD dijawab tepat sesuai jenis bahaya pada kasus. & Langkah pengendalian umumnya benar, tetapi pemilihan APD belum sesuai jenis bahaya. & Langkah pengendalian atau pemilihan APD keliru. \\ \\
Ketepatan konsep organisasi dan komunikasi K3 & 4 & Peran organisasi, komunikasi, dan personel K3 menuju zero accident dijelaskan tepat dan lengkap. & Peran organisasi dan komunikasi K3 cukup benar, tetapi penjelasan belum lengkap. & Peran organisasi atau komunikasi K3 dijelaskan keliru. \\}

\assessmentblock{Case Method Perbaikan SMK3 dan K3 Industri 4.0}{Case Method}{SCPMK0101-00804}{Mahasiswa menganalisis kasus perbaikan sistem manajemen K3 dan analisis biaya pada sebuah industri, lalu menyusun usulan implementasi K3 yang memanfaatkan teknologi otomatisasi dan pertukaran data era industri 4.0.}{Menganalisis kasus secara berkelompok, mengidentifikasi kelemahan SMK3, menghitung komponen biaya, menyusun usulan perbaikan berbasis teknologi industri 4.0, dan mempresentasikan hasil.}{Laporan analisis kasus dan bahan presentasi.}{Ketepatan analisis kasus SMK3 dan analisis biaya & 4 & Kelemahan SMK3 diidentifikasi berbasis data kasus dan komponen biaya K3 dihitung dengan asumsi yang jelas. & Kelemahan SMK3 teridentifikasi, tetapi analisis biaya atau asumsinya belum lengkap. & Analisis tidak berbasis data kasus atau komponen biaya tidak dihitung. \\ \\
Kelayakan usulan perbaikan dan implementasi industri 4.0 & 3 & Usulan perbaikan operasional, memanfaatkan teknologi otomatisasi/pertukaran data, dan sesuai regulasi K3. & Usulan cukup relevan, tetapi pemanfaatan teknologi atau kesesuaian regulasi belum jelas. & Usulan tidak operasional atau tidak terkait teknologi industri 4.0. \\ \\
Kualitas presentasi dan kerja kelompok & 3 & Presentasi runtut, argumentasi didukung data, dan kontribusi anggota kelompok merata. & Presentasi cukup jelas, tetapi argumentasi atau pembagian kontribusi belum merata. & Presentasi tidak menjelaskan hasil analisis atau kontribusi tidak terlihat. \\}

\assessmentblock{UAS Penerapan Pengendalian dan Rancangan Program K3}{UAS (ujian tulis daring)}{SCPMK0101-00803, SCPMK0101-00804}{Mahasiswa mengerjakan ujian akhir semester yang mencakup penerapan manajemen risiko, APD, protokol kesehatan, metode statistik pengendalian K3, serta perancangan program kerja, organisasi, dan audit K3.}{Mengerjakan ujian tulis daring secara mandiri, menjawab soal penerapan dan soal perancangan berbasis kasus.}{Lembar jawaban ujian daring.}{Ketepatan penerapan pengendalian risiko, APD, dan protokol kesehatan & 10 & Langkah pengendalian risiko, pemilihan APD, dan protokol kesehatan diterapkan tepat sesuai kasus dan regulasi. & Penerapan umumnya benar, tetapi sebagian langkah tidak sesuai kasus atau regulasi. & Penerapan pengendalian keliru atau mengabaikan regulasi K3. \\ \\
Kualitas rancangan program kerja, organisasi, dan audit K3 & 10 & Program kerja, struktur organisasi, tupoksi, dan rencana audit K3 dirancang lengkap, runtut, dan dapat diterapkan. & Rancangan tersedia, tetapi kelengkapan tupoksi atau rencana audit masih kurang. & Rancangan tidak operasional atau komponen utama tidak ada. \\ \\
Kualitas jawaban kasus terintegrasi K3 industri 4.0 & 10 & Solusi kasus mengintegrasikan pengendalian K3, aspek biaya, dan teknologi industri 4.0 secara konsisten. & Solusi kasus cukup relevan, tetapi integrasi antar aspek belum konsisten. & Solusi kasus tidak relevan atau hanya menyalin teori. \\}

\begin{tblr}{width=\textwidth,colspec={|Q[l,m,3.2cm]|X[l,m]|},hlines,vlines,hspan=minimal,row{1-3}={bg=headgray},rowsep=1.5pt,colsep=3pt}
Jadwal Pelaksanaan: & Minggu 1--17 sesuai RPS. Setiap evaluasi memiliki RTM dan rubrik multi-kriteria. \\
Lain-Lain Yang Diperlukan: & LMS, alat peraga keselamatan kerja, peraturan perundang-undangan K3, dan perangkat presentasi. \\
Pustaka: & Mengacu pustaka utama dan pendukung pada bagian RPS. \\
\end{tblr}
